\documentclass{article}

\textwidth=6.5in
\textheight=8.5in
\voffset=-54pt
\oddsidemargin=0in
\evensidemargin=0in
\setlength{\parskip}{8pt}
\setlength{\parindent}{0pt}

\usepackage{workbook}

\usepackage[workshop]{handout}

\begin{document}

\htitle{Workshop 1: Related Rates of Change and Implicit Differentiation}
\begin{workshop}
    The objective for this week's workshop is to have students practice solving related rates problems. We have seen mostly examples of related rates problems that make use of the geometry of right triangles, similar triangles, and cones, which are the three types of related rates problems students can practice in workshop today. There is also a different type of implicit differentiation problem at the end of this worksheet, but nonetheless students will practice this overall skill with related rates problems as well. Expect students to have different levels of comfort and fluency with related rates. To smooth the transition into problem solving, \#1 contains scaffolds that walk students through some key steps that they can follow in \#2 and \#3, for example. 
    \newline
    \newline
    Here is an example lesson plan for workshop:
    \begin{itemize}
        \item (up to 10 minutes) Set the tone for problem-solving and norms for the remainder of the semester, including attendance (mandatory), timeliness, and collaboration. Take attendance, and give a general overview of the overarching goal of related rates problems: \textit{``In related rates problems, we typically look at scenarios where multiple quantities change with respect to time, and we are interested in \underline{relating} their \underline{rates} of change, which is where the phrase ``related rates" comes from! In class and on your p-set, you started these problems by naming variables for these quantities that change, and then used things like the Pythagorean Theorem, similar triangles, and other formulas to relate these variables. In order to relate their \textit{rates of change}, we use implicit differentiation to take the derivative of both sides of the equation that relates these variables, usually with respect to time, and then solve for our goal using the information we were given in the problem. Today we're going to practice applying this overall strategy."}
        \item (30 minutes) Give students about 15 minutes to work on \#1 at the boards, and circulate to provide group support. Collect the class after most groups seem to be finishing up \#1, and then thoroughly debrief the problem solving process using input from student work and discussion, prioritizing a clean and organized model of board work that students can reference throughout workshop.
        \item (30 minutes) Prompt students to work on \#2 in groups using the problem solving strategies from \#1. If you notice there are some components of \#2 that would benefit from a whole-class debrief, feel free to collect the class for a moment and address them. However, students can work at their own pace on the remainder of the worksheet as you circulate and discuss work with students. 
        \item After completing \#2, there is a prompt to have students reflect on whether they would like to solve an additional related rates problem first, or tackle a different type of implicit differentiation problem in \#4. 
    \end{itemize}
\end{workshop}
\begin{enumerate}
\item A coffee bean reservoir takes the shape of an inverted cone. The radius at the top of the reservoir is $4$ centimeters and the height is $20$ centimeters. If coffee beans enter the reservoir at a rate of $2$ cubic centimeters per second, how fast is the height of the beans changing when the height is $10$ centimeters?

\underline{Note}: The volume of a cone is $V=\frac{1}{3}\pi r^2h$.
\begin{enumerate}
\item First, draw a general picture of a partially-filled coffee bean reservoir. Label the dimensions that remain constant and the dimensions that vary with respect to time. 
\newline
\probonly{
\begin{center}
\includegraphics[width=2in]{images/spring2012-midterm1-meghan-cone.pdf}
\end{center}
}

\begin{solution}
    \begin{center}
        \begin{tikzpicture}[x=0.3cm, y=0.3cm]
          \def\R{8}
          \def\H{12}
          \def\r{5}
          \def\h{60/8}
          \def\aspect{4}
          \draw [fill=white] (-\R,\H) -- (\R,\H) -- (0,0) -- cycle;
          \draw [fill=white] (0, \H) circle [x radius=\R, y
          radius={\R/\aspect}];
          \draw [fill=brown!40!white,opacity=1] (-\r, \h) -- (\r, \h) --
          (0,0) -- cycle; 
          \draw [fill=brown!20!white,opacity=1, draw opacity=0] (0, \h) ellipse
          [x radius=\r, 
          y radius={\r/\aspect}];
          \draw (-\r, \h) arc [x radius=\r, y radius={\r/\aspect}, start
          angle=180, end angle=360];      
          \draw[dashed] (\r, \h) arc [x radius=\r, y radius={\r/\aspect},
          start angle=0, end angle=180];
          \draw (0, \H) -- node[above, inner sep=0.1pt]{4 cm} (\R, \H);
          \draw[|-|, xshift=-12pt] (-\R, 0) -- node[left]{20 cm} (-\R,
          \H);
          \draw[red] (0, \h) -- node[above, inner sep=0]{$r$} (\r, \h);
          \draw[red, |-|, xshift=-12pt] (-\r, 0) -- node[left]{$h$} (-\r,
          \h);
          \draw (0, 0) node[below, font=\normalsize]{\emph{Generic}};
        \end{tikzpicture} \hspace{0.5in}
        \end{center}
\end{solution}
\begin{workshop}
    This question is a good checkpoint to remind students about derivative notation and how to use it to streamline the problem solving process. For example: ``I want to know $\frac{dh}{dt}$ when $h=10$cm, and I know that $\frac{dV}{dt}=2$ cm$^3$/sec." 
\end{workshop}
\item Using the variable names you established above, identify your goal and a rate of change you already know with correct derivative notation and units.

\begin{solution}
    Our goal is to find $\dfrac{dh}{dt}$ when $h=10$ cm. We know that $\dfrac{dv}{dt}=2$ cm$^3$/sec.
    \end{solution}
\pspace{\vfill\vfill}
\item Draw a snapshot of the reservoir at the moment we're interested in. Remember to use constants to  label the dimensions we know that are true only at that one moment.
\probonly{\begin{center}
    

\includegraphics[width=2in]{images/spring2012-midterm1-meghan-cone.pdf}
\end{center}
}

\begin{solution}
    \begin{center}
        \begin{tikzpicture}[x=0.3cm, y=0.3cm]
          \def\R{8}
          \def\H{12}
          \def\r{5}
          \def\h{60/8}
          \def\aspect{4}
          \draw [fill=white] (-\R,\H) -- (\R,\H) -- (0,0) -- cycle;
          \draw [fill=white] (0, \H) circle [x radius=\R, y
          radius={\R/\aspect}];
          \draw [fill=brown!40!white,opacity=1] (-\r, \h) -- (\r, \h) --
          (0,0) -- cycle; 
          \draw [fill=brown!20!white,opacity=1, draw opacity=0] (0, \h) ellipse
          [x radius=\r, 
          y radius={\r/\aspect}];
          \draw (-\r, \h) arc [x radius=\r, y radius={\r/\aspect}, start
          angle=180, end angle=360];      
          \draw[dashed] (\r, \h) arc [x radius=\r, y radius={\r/\aspect},
          start angle=0, end angle=180];
          \draw (0, \H) -- node[above, inner sep=0.1pt]{4 cm} (\R, \H);
          \draw[|-|, xshift=-12pt] (-\R, 0) -- node[left]{20 cm} (-\R,
          \H);
          \draw[red] (0, \h) -- node[above, inner sep=0]{$2$ cm} (\r, \h);
          \draw[red, |-|, xshift=-12pt] (-\r, 0) -- node[left]{$10$cm} (-\r,
          \h);
          \draw (0, 0) node[below, font=\normalsize]{\emph{Snapshot}};
        \end{tikzpicture} \hspace{0.5in}
        \end{center}
\end{solution}

\pspace{\newpage}

\item Which two variables will you have to relate to calculate your goal, given the rate you already know? Write an equation that relates these two variables. (How are similar triangles helpful here?)

\begin{solution}
    \begin{itemize}
    \item We want to relate $V$, the volume of the cone, with its height $h$. Note that the formula $V=\frac{1}{3}\pi r^2h$ expresses $V$ in terms of both $r$ and $h$; here, let's plan to substitute an expression for $r$ in terms of $h$ using the properties of similar triangles within the cone.
    \item Because the ratio of the height to the radius of this cone is constant, we know that $\frac{r}{h}=\frac{4}{20}=\frac{1}{5}$, so $r=\frac{h}{5}$.
    \item Let's substitute $\frac{h}{5}$ for $r$ in the volume formula $V=\frac{1}{3}\pi r^2h$ to yield $V=\frac{1}{3}\pi \left(\frac{h}{5}\right)^2h$, which simplifies to $V=\frac{1}{3}\pi \left(\frac{h^2}{25}\right)h=\frac{1}{75}\pi h^3$.
    \end{itemize}
\end{solution}
\probonly{
\vfill}

\item Use implicit differentiation on the equation from part (d) and solve for your goal.

\begin{solution}


 Now that we have a formula for $V$ in terms of $h$, let's implicitly differentiate with respect to time to generate an equation that contains what we know, $\dfrac{dV}{dt}$, and what we want to find, $\dfrac{dh}{dt}$.
\begin{align*}
    \frac{d}{dt} V&= \dfrac{d}{dt} \left(\frac{1}{75}\pi h^3\right) \\
    \dfrac{dV}{dt}&=\frac{1}{25}\pi h^2 \frac{dh}{dt}
\end{align*}
Solving for $\dfrac{dh}{dt}$ when $h=10$, 
\begin{align*}
    \frac{dh}{dt}&=\dfrac{\frac{dV}{dt}}{\frac{1}{25}\pi h^2} \\[4pt]
    &=\dfrac{2}{\frac{1}{25}\pi (10)^2} \\[4pt]
    &=\dfrac{2}{4\pi} \\[4pt]
    &=\dfrac{1}{2\pi} \text{ cm/sec}\\[4pt]
\end{align*}
    
\end{solution}
\probonly{
\vfill} 
\end{enumerate}
\pspace{\newpage}



\item Mike stands 12 feet behind the pick up spot for his iced coffee. The coffee bar, where a barista makes Mike's drink, is perpendicular to Mike's path towards the pick up spot. At the moment Mike begins to shuffle towards the pick up spot at a rate of $\frac{1}{2}$ feet per second, the barista rushes Mike's drink to the pick up spot from 5 feet away at a rate of 2 feet per second. At this moment, how quickly is the distance between Mike and his coffee changing?

\underline{Note:} Be sure to assign and define variables for quantities that change with respect to time, include a generic diagram and a snapshot, and lay out your work in a clear, logical, and step-by-step manner.

\begin{solution}
\begin{itemize}
\item Let's assign variables $x$, $y$, and $z$ to the distances in this scenario that measure the three sides of the right triangle set-up of the coffee pick-up. While we're at it, let's also fill in the corresponding side lengths of this right triangle at the moment we are interested in relating rates of change: 

\begin{center}
    \begin{multicols}{2}
\begin{tikzpicture}[scale=0.3]
     \draw[thick](0,0) -- (0,12) -- (5,12) -- (0,0);
     \node[below] at (0,0){Mike};
     \node[left] at (0,12){pick-up spot};
     \node[right] at (5,12){barista};
     \draw[fill] (0,0) circle [radius=4pt];
     \draw[fill] (0,12) circle [radius=4pt];
     \draw[fill] (5,12) circle [radius=4pt];
     \node[left] at (0,6){$x$};
     \node[above] at (2.5,12){$y$};
     \node[right] at (2.5,6){$z$};
     \draw (0, -2) node[below, font=\normalsize]{\emph{Generic}};
     \end{tikzpicture}
 \begin{tikzpicture}[scale=0.3]
     \draw[thick](0,0) -- (0,12) -- (5,12) -- (0,0);
     \node[below] at (0,0){Mike};
     \node[left] at (0,12){pick-up spot};
     \node[right] at (5,12){barista};
     \draw[fill] (0,0) circle [radius=4pt];
     \draw[fill] (0,12) circle [radius=4pt];
     \draw[fill] (5,12) circle [radius=4pt];
     \node[left] at (0,6){$12$ ft};
     \node[above] at (2.5,12){$5$ ft};
     \node[right] at (2.5,6){$13$ ft};
     \draw (0, -2) node[below, font=\normalsize]{\emph{Snapshot}};
 \end{tikzpicture}
 \end{multicols}
 \end{center}
 \item According to the variable names in our generic picture, we want to find $\frac{dz}{dt}$ when $x=12$, $y=5$, and $z=13$. At this moment, we also know that $\frac{dx}{dt}=-\frac{1}{2}$ ft/sec and $\frac{dy}{dt}=-2$ ft/sec. (Note that both of these rates are negative because the overall distances that these rates of change trace are \textit{decreasing} over time.)
 \item The variables $x$, $y$, and $z$ are related by the Pythagorean Theorem, where $x^2+y^2=z^2$. Let's differentiate both sides with respect to time to generate an equation that contains our goal and the rates we already know, plug in these known rates and lengths, and solve for $\frac{dz}{dt}$:
 \begin{align*}
     \dfrac{d}{dt} \left(x^2+y^2\right)&=\dfrac{d}{dt} z^2 \\
     2x\frac{dx}{dt}+2y\frac{dy}{dt}&=2z\frac{dz}{dt} \\
      x\frac{dx}{dt}+y\frac{dy}{dt}&=z\frac{dz}{dt} \\
     (12)\left(-\frac{1}{2}\right)+(5)(-2)&=(13)\frac{dz}{dt} \\
     \frac{dz}{dt}&=-\frac{16}{13} \text{ft/sec} 
 \end{align*}
 \end{itemize}
\end{solution}
\pspace{\vfill}
\probonly{
\begin{framed}
    \textbf{Reflect} If you'd prefer to do more practice with modeling and related rates of change, proceed to \#3 first. If you would like to practice computation with implicit differentiation, proceed to \#4 first.
\end{framed}}
\newpage

% from \ProblemBank23\RelatedRates(NoTrig).tex

\item 
On a dark night, a glowing UFO is hovering $30$ feet over a field.  A $6$-foot-tall farmer comes running to see what is happening.  When the farmer is $60$ feet away from the spot on the ground directly under the UFO, he is running at $8$ feet per second toward the UFO.

\begin{center}
\includegraphics[width=3in]{images/ufo-running.pdf}
\end{center}
\begin{enumerate}
\item  \begin{problem}[2in] How long is the farmer's shadow (cast by the light from the UFO) at this moment?\end{problem}

\begin{solution}
Let $s$ represent the length of the shadow. Then we can use similar triangles: 
\begin{align*}
\frac6s &= \frac{30}{60+s}\\
360+6s&=30s\\
360&=24s\\
s&=15
\end{align*}
The farmer's shadow is 15 feet long.
\end{solution}

\item  \begin{problem}[2in]  How quickly is the length of the shadow changing at this moment?  Be sure to specify whether the shadow is getting longer or shorter.

\underline{Note:} Assign and define variables for quantities that change with respect to time, include a generic diagram and a snapshot, and lay out your work in a clear, logical, and step-by-step manner.
\end{problem}

\begin{solution}
Let $x$ be the horizontal distance between the man and the spot on the ground beneath the UFO. Let $s$ be the length of his shadow. Then we know that $\frac{dx}{dt} = -8$ feet per second and we want to know what $\frac{ds}{dt}$ is.

Again, we can use similar triangles: 
\begin{align*}
\frac6s &= \frac{30}{x+s}\\
6x+6s&=30s\\
6x&=24s\\
\end{align*}
We differentiate both sides with respect to $t$:
\begin{align*}
\frac{dx}{dt}&=4\frac{ds}{dt}\\
-8&=4\frac{ds}{dt}\\
\frac{ds}{dt} &=-2
\end{align*}
Therefore, the shadow is getting shorter at a rate of $2$ feet per second.
\end{solution}

\begin{center}
\includegraphics[width=3in]{images/ufo-standing.pdf}
\end{center}
\begin{workshop}
    Students might be hesitant to compute $\frac{d}{dt}sh$ (once it appears in their work) as $h\frac{ds}{dt}+s\frac{dh}{dt}$ because they may not have encountered this product rule computation in the related rates context thus far; however, you can remind students that this is the Chain Rule in action, and that both variables and their rates of change are relevant to solve for the goal.
\end{workshop}
\item  \begin{problem} \textit{Challenge} When the 6-foot-tall farmer is $30$ feet away from the spot on the ground directly under the UFO, he becomes frightened and stops where he is. The UFO then starts descending towards Earth at a rate of $2$ feet per second.  When the UFO is $24$ feet over the field, the farmer's shadow is $10$ feet long. At this moment, how quickly is the farmer's shadow changing in length?  Be sure to specify whether the shadow is getting longer or shorter.

\underline{Note:} Assign and define variables for quantities that change with respect to time, include a generic diagram and a snapshot, and lay out your work in a clear, logical, and step-by-step manner.
\end{problem}

\begin{solution} 
Again let $s$ represent the length of the shadow, and this time let $h$ represent the height of the UFO. Then we know that $\frac{dh}{dt} = -2$ and we want to know $\frac{ds}{dt}$. Using similar triangles again we have: 
\begin{align*}
\frac{6}{s} &= \frac{h}{30+s} \\
 180+6s &= sh\\
6\frac{ds}{dt} &= h\frac{ds}{dt} + s\frac{dh}{dt}\\
6\frac{ds}{dt}&= 24\frac{ds}{dt} -20\\
20 &=18 \frac{ds}{dt}\\
\frac{ds}{dt} &=\frac{10}9
\end{align*}
Therefore, the shadow is now getting longer at a rate of $\frac{10}{9}$ feet per second.
\end{solution}
\end{enumerate}


\pspace{\newpage}
\item Let's look at the implicit curve for the lima bean we saw in class, defined by $(x^2 + y^2)^2 = x^3 + y^3$.
\begin{center}
\probonly{
    \includegraphics[scale=0.75]{images/implicit-curve-crooked-egg.pdf}}
    

\end{center}
\begin{center}
\solonly{
    \includegraphics[scale=0.2]{images/tanglima2.png}}
    

\end{center}

\begin{enumerate} 
\item Draw a line tangent to the lima bean implicit curve at $\left(\dfrac{1}{2}, \dfrac{1}{2}\right)$.
\item Find the equation of the line you plotted in part (a). 

\begin{solution}
    Let's use implicit differentiation to take the derivative of $(x^2 + y^2)^2 = x^3 + y^3$ with respect to $x$, then plug in $x=\frac{1}{2}$ and $y=\frac{1}{2}$ and simplify:
    \begin{align*}
        \frac{d}{dx}[(x^2 + y^2)^2] &= \frac{d}{dx}[x^3 + y^3] \\
        2(x^2+y^2)(2x+2y\cdot\frac{dy}{dx}) &= 3x^2+3y^2\cdot\frac{dy}{dx}\\ 
        2\left(\left(\frac{1}{2}\right)^2+\left(\frac{1}{2}\right)^2\right)\left(2\left(\frac{1}{2}\right)+2\left(\frac{1}{2}\right)\cdot\frac{dy}{dx}\right) &=3\left(\frac{1}{2}\right)^2+3\left(\frac{1}{2}\right)^2\cdot\frac{dy}{dx} \\
        (1)\left(1+\frac{dy}{dx}\right)&=\frac{3}{4}+\frac{3}{4}\frac{dy}{dx}
    \end{align*}
    Now, let's collect like terms and solve for $\frac{dy}{dx}$:
    \begin{align*}
          1+\frac{dy}{dx}&=\frac{3}{4}+\frac{3}{4}\frac{dy}{dx}\\
          -\frac{1}{4}&=\frac{1}{4}\frac{dy}{dx}\\
          \frac{dy}{dx}&=-1
    \end{align*}
The equation of the line with a slope of $-1$ that passes through $\left(\dfrac{1}{2}, \dfrac{1}{2}\right)$ is $y-\dfrac{1}{2}=-1\left(x-\dfrac{1}{2}\right)$, or $T(x)=-x+1$.
\end{solution}
\pspace{\vfill \vfill \vfill }
\item A bug walking down the lima bean passes through the point $\left(\dfrac{1}{2}, \dfrac{1}{2}\right)$. Use linear approximation to estimate its $y$-coordinate when $x=\dfrac{3}{5}$. (How is part (b) useful?)

\begin{solution} We can use the tangent line $T(x)=-x+1$ as an approximation for the lima bean curve near the point $\left(\dfrac{1}{2}, \dfrac{1}{2}\right)$. Thus, when the bug's $x$-coordinate on the lima bean curve is $x=\frac{3}{5}$, its $y$- coordinate is approximately $T\left(\frac{3}{5}\right)=-\frac{3}{5}+1=\frac{2}{5}$.
\end{solution}

\pspace{\vfill}

\end{enumerate}
\end{enumerate}
\end{document}
